\documentclass[12pt]{article}
\usepackage[utf8]{inputenc}
\usepackage[T1]{fontenc}
\usepackage[margin=1in]{geometry}
\usepackage{setspace}
\usepackage{xcolor}
\usepackage[colorlinks=true, urlcolor=blue!60!black, linkcolor=blue!60!black]{hyperref}
\setlength{\parindent}{2em}
\setlength{\parskip}{1.25em}
\renewcommand{\baselinestretch}{1.0}
\begin{document}
\begin{center}
{\LARGE \textit{Crosby v. National Foreign Trade Council}, 530 U.S.\ 363 (2000)\par}
\vspace{0.8em}
{\large Case Brief --- Four-Page Edition\par}
\end{center}
\vspace{1em}

\section{Memory Jogger}

A unanimous Court held that Massachusetts's law restricting state purchases from companies doing business with Burma was preempted because it obstructed the federal Burma Act's grant of flexible presidential discretion, its calibrated sanctions, and its call for multilateral diplomacy (\href{https://supreme.justia.com/cases/federal/us/530/363/}{Justia}).

\section{Detailed Case Facts}

In June 1996, reacting to human rights abuses by Burma's military junta, Massachusetts barred state agencies from buying from companies on a ``restricted purchase list'' of entities ``doing business with Burma,'' modeled on anti-apartheid selective-purchasing laws. Listed companies were effectively disqualified unless they bid at least 10\% below the lowest unlisted bidder. The law covered domestic and foreign firms, parents and subsidiaries, and pre-existing investments, and had no sunset or waiver.

Three months later Congress enacted \S~570 of the 1997 Foreign Operations Appropriations Act. It imposed immediate limited sanctions; authorized the President to ban ``new investment'' by ``United States persons'' upon certifying repression; directed a ``comprehensive, multilateral strategy'' with ASEAN and others; and let the President waive sanctions ``in the national security interests of the United States.'' President Clinton banned new investment by executive order in 1997. The European Union and Japan protested the Massachusetts law and initiated WTO proceedings.

\section{Procedural History}

The National Foreign Trade Council, 34 of whose members were on the restricted list, sued, invoking the foreign affairs power under \textit{Zschernig v.\ Miller}, the dormant Foreign Commerce Clause, and preemption. The district court enjoined the law, and the First Circuit affirmed on all three grounds. The Supreme Court affirmed on June 19, 2000, solely on statutory preemption.

\section{Judicial Votes}

Affirmed, \textbf{9--0}.
\begin{itemize}
\item \textbf{Majority:} Justice Souter, joined by Chief Justice Rehnquist and Justices Stevens, O'Connor, Kennedy, Ginsburg, and Breyer.
\item \textbf{Concurring in the judgment:} Justice Scalia, joined by Justice Thomas.
\end{itemize}

\section{Holding}

The Massachusetts Burma Law is preempted because it stands as an obstacle to the full purposes and objectives of the federal Act: its delegation of flexible discretion to the President, its limitation of sanctions to U.S.\ persons and new investment, and its directive to pursue multilateral diplomacy. The Court reserved whether a presumption against preemption applies in foreign affairs, field preemption, the \textit{Zschernig} theory, and the dormant Foreign Commerce Clause.

\section{Analysis of Opinions}

\textbf{Majority (Souter).} A state law must yield when it obstructs the objectives of Congress (\textit{Hines v.\ Davidowitz}), judged by the statute's purpose and ``natural effect.'' Under \textit{Wisconsin v.\ Gould}, using purchasing power instead of regulation is no defense. First, invoking Justice Jackson's \textit{Youngstown} framework, the President acting with congressional authorization is at his maximum; the rigid, perpetual state law, in \textit{Dames \& Moore}'s metaphor, reduces the value of the President's ``bargaining chips.'' Second, the state law reached pre-existing investments, ordinary trade, and foreign firms that Congress excluded, and a shared goal does not cure conflicting means. Third, foreign protests and Executive representations were ``competent and direct evidence'' that the law impaired the President's ability to ``speak for the United States among the world's nations''; \textit{Barclays Bank} was distinguished because Congress had not acted there. Congressional silence about the state law was ``ambiguous,'' not permission.

\textbf{Scalia, with Thomas (concurring in the judgment).} The conclusion was ``perfectly obvious on the face of this statute,'' so the dozen footnotes of legislative history were illegitimate and unnecessary, like St.\ Augustine stealing pears ``just for the devil of it,'' and they inflate litigation costs.

\textbf{Later cases.} \textit{Garamendi} (2003) extended the reasoning to executive agreements; \textit{Arizona v.\ United States} (2012) used its obstacle framework; \textit{Zivotofsky v.\ Kerry} (2015) quoted its ``one voice'' principle.

\section{Examples: Future Applications}

\begin{enumerate}
\item \textbf{Same-Side: Broader, Unwaivable Ban.} New York bars contracts with any firm doing business in a country Congress sanctioned narrowly with a waiver. The law blunts leverage, overrides calibration, and impedes diplomacy.
\item \textbf{Same-Side: Matching Scope, No Waiver.} A State copies a federal investment ban but omits the waiver; when the President waives the federal ban, the state law is preempted to the extent it outlasts the waiver.
\item \textbf{Opposite-Side: Express Authorization.} Congress expressly authorizes matching state pension divestment that sunsets on federal delisting. It survives.
\item \textbf{Opposite-Side: No Federal Statute.} A State boycotts a country Congress has not sanctioned. \textit{Crosby} does not apply; challengers must press the reserved \textit{Zschernig} and Foreign Commerce theories.
\item \textbf{Fence-Sitter: Tracking the Federal List.} A State divests from firms on a Treasury sanctions list and reinvests upon delisting. It preserves bargaining chips, but adds pressure Congress did not impose, and silence is not permission.
\end{enumerate}

\section{Critique}

\textbf{Market-participant asymmetry.} States enjoy a ``market participant'' shield under the dormant Commerce Clause, yet \textit{Crosby}, following \textit{Gould}, confirms that spending-power choices enjoy no immunity from Supremacy Clause preemption.

\textbf{Elasticity and foreign protests.} Treating any partial federal sanction as a ``calibration of force'' invites expansive obstacle preemption, and crediting foreign diplomatic protests risks giving foreign sovereigns a de facto veto over subnational policy.

\section{Key Quotations}

\begin{quote}
``The fact of a common end hardly neutralizes conflicting means, see \textit{Gade v.\ National Solid Wastes Management Assn.}, 505 U.S.\ 88, 103 (1992), and the fact that some companies may be able to comply with both sets of sanctions does not mean that the state Act is not at odds with achievement of the federal decision about the right degree of pressure to employ. \dots\ Sanctions are drawn not only to bar what they prohibit but to allow what they permit, and the inconsistency of sanctions here undermines the congressional calibration of force.''
\end{quote}
\noindent --- Justice Souter, Opinion of the Court, 530 U.S.\ at 379--380.

\begin{quote}
``We find that the state law undermines the intended purpose and `natural effect' of at least three provisions of the federal Act, that is, its delegation of effective discretion to the President to control economic sanctions against Burma, its limitation of sanctions solely to United States persons and new investment, and its directive to the President to proceed diplomatically in developing a comprehensive, multilateral strategy toward Burma.''
\end{quote}
\noindent --- Justice Souter, Opinion of the Court, 530 U.S.\ at 373--374.

\begin{quote}
``A failure to provide for preemption expressly may reflect nothing more than the settled character of implied preemption doctrine that courts will dependably apply, and in any event, the existence of conflict cognizable under the Supremacy Clause does not depend on express congressional recognition that federal and state law may conflict, \textit{Hines}, 312 U.S., at 67. The State's inference of congressional intent is unwarranted here, therefore, simply because the silence of Congress is ambiguous.''
\end{quote}
\noindent --- Justice Souter, Opinion of the Court, 530 U.S.\ at 387--388.

\end{document}
